\section{Comisión de revisión de protocolos: visión global}
Este manual reúne lo que el actor hace en \textbf{todas} las opciones de titulación y señala, al final, lo particular de cada una. Para el detalle paso a paso de una opción, consulta el manual específico de esa opción.

\subsection{Opciones en las que participa}
\begin{itemize}
  \item Artículo especializado para publicar en revista indizada
  \item Ensayo
  \item Memoria de experiencia laboral
  \item Reporte de aplicación de conocimientos
  \item Reporte de autoempleo profesional
  \item Reporte de residencia de investigación
  \item Tesina
  \item Tesis
\end{itemize}

\subsection{Pasos comunes}
Recibe el protocolo turnado por Titulación, asigna evaluadores con fecha límite (siete días hábiles por defecto), da seguimiento con el semáforo, y consolida y firma el dictamen cuando todos firmaron.

\subsection{Particularidades por opción}
Participa en las opciones con trabajo escrito; el flujo es el mismo en todas, con la diferencia del tipo de trabajo y de las evidencias propias de cada una:
\subsubsection*{Artículo especializado para publicar en revista indizada}
Artículo original e inédito, enviado a una revista inscrita en índices validados.
\subsubsection*{Ensayo}
Exposición crítica y propositiva de un tema, con postura intelectual propia.
\subsubsection*{Memoria de experiencia laboral}
Recopilación y resultados de una experiencia profesional en un ámbito reconocido.
\subsubsection*{Reporte de aplicación de conocimientos}
Propuesta de solución a un problema determinado, derivada de lo aprendido en la carrera.
\subsubsection*{Reporte de autoempleo profesional}
Informe sobre la constitución y desarrollo de una unidad económica formal.
\subsubsection*{Reporte de residencia de investigación}
Informe de actividades de apoyo a la investigación bajo la dirección de un investigador activo en el SNI.
\subsubsection*{Tesina}
Investigación documental de carácter monográfico, con aportación al estado del arte.
\subsubsection*{Tesis}
Investigación documental, experimental, empírica o teórica; individual o de hasta tres pasantes.
